%% ma: L12-B08-0040
%% lop: 12
%% chuong: 2
%% bai: 8
%% dang: 12.8.8
%% loai: TN4
%% muc: TH
%% dap_an: "A"
%% nguon: "LaTeX giáo viên gửi 21/10/2026 (M05, Vectơ và hệ trục tọa độ trong không gian, 107 câu), Phần I Câu 31"
%% kien_thuc: "Bài 6 (quy tắc ba điểm), Bài 8 (biểu thức tọa độ các phép toán vectơ)"
%% trang_thai: chinh-thuc
%% ghi_chu: "Trùng ý tưởng với L12-B08-0010 (hệ thức MA + kMB = 0), mở rộng cho ba điểm, nên nhập cùng dạng 12.8.8. Phần I Câu 12 của tài liệu (2IA + IB = 0) chỉ đổi số so với L12-B08-0010 nên đã bỏ"
\begin{ex}
Trong không gian $Oxyz$, cho ba điểm $A(1;0;0)$, $B(0;2;0)$, $C(0;0;3)$. Gọi $I$ là điểm thỏa mãn hệ thức $\overrightarrow{IA}+2\overrightarrow{IB}+3\overrightarrow{IC}=\overrightarrow{0}$. Tọa độ điểm $I$ là
\choice{\True $\left(\dfrac16;\dfrac23;\dfrac32\right)$}{$(1;4;9)$}{$\left(\dfrac16;\dfrac13;\dfrac12\right)$}{$\left(\dfrac13;\dfrac23;1\right)$}
\loigiai{$\overrightarrow{IA}+2\overrightarrow{IB}+3\overrightarrow{IC}=\overrightarrow{0}\Leftrightarrow 6\overrightarrow{OI}=\overrightarrow{OA}+2\overrightarrow{OB}+3\overrightarrow{OC}=(1;4;9)$. Vậy $I\left(\dfrac16;\dfrac23;\dfrac32\right)$.}
\end{ex}
